\scanpage{79}
\mspage{28}

\uline{Définition 2} Sous les conditions de la prop. 3, on dit que $f$ est $W$-\uline{propre}. On dit que $f$ est $W$-lisse si $f^\circ$ est $W$-propre, i.e. si les morphismes d'\unsure{inclusions|inclusion} (\unsure{*|x}) sont \uline{$W$-asphériques} etc (conditions duales de celles de la prop 3).

\uline{Remarque} La condition a) de la prop 3 est satisfaite en particulier si l'inclusion \inserted{(*)} a un adjoint à gauche (cf remarque page 23). Cela signifie que les catégories $X(x_0,u_0)$ ont un \unsure{objet} initial. Cela signifie aussi que la catégorie $X$ sur $Y$ est $\mathrm{Cat}$-cofibrée sur $Y$ au sens faible, i.e. qu'il existe, pour $u_0:y_0\mathrel{\relbar}y_1$ flèche dans $Y$, \deleted{un foncteur} et \unsure{$a_0$|$x_0$} dans $X_{y_0}$, un objet $u_{0*}(x_0)$ dans $X_{y_1}$ qui représente le foncteur covariant $X_{y_1}\mathrel{\relbar}(\mathrm{Ens})$, $x_1\longrightarrow\mathrm{Hom}_{\unsure{u|u_0}}(x_0,x_1)$. Ainsi, si $X$ est $\mathrm{Cat}$-cofibrée au sens faible sur $Y$, et a fortiori si elle est $\mathrm{Cat}$-cofibrée, alors $f:X\longrightarrow Y$ est $W$-propre. Dualement, si $f$ est faiblement $\mathrm{Cat}$-fibrée, $f$ est lisse.

% slanted
\margin{\unsure{Dire aussi|D'où, aussi|D'où, on voit} : Si $X\xrightarrow{f}Y\xrightarrow{g}Z$, si $gf$ et $f$ $W$-lisses, alors $g$ est lisse (en les points de $\mathrm{Im}\,f$ ?), idem pour propre.} \uline{Proposition 4} \deleted{Si} La notion de morphisme \deleted{\unsure{fa}}propre est stable par composition, par changement de base. Les coisom locaux sont propres, \inserted{(\unsure{trivial dans b|trivial selon b})} a fortiori les isom, sont propres, les imm. fermées sont propres.

Enoncé dual pour les morphismes \unsure{lisse|lisses} : \unsure{stabilité} par chgt de base, par composition, les iso. locaux sont lisses, notamment les imm. ouvertes.

% archive sheet number: 79
